% Cross-Lingual Factual Consistency of LLMs: A Three-Way Script Study on Hindi
% Threws Odyssey 2026 — Applied Research Submission

\documentclass[11pt,a4paper]{article}
\usepackage[utf8]{inputenc}
\usepackage[T1]{fontenc}
\usepackage{amsmath,amssymb}
\usepackage{graphicx}
\usepackage{booktabs}
\usepackage{hyperref}
\usepackage{xcolor}
\usepackage{multirow}
\usepackage{subcaption}
\usepackage[margin=1in]{geometry}
\usepackage{natbib}
\bibliographystyle{plainnat}

\title{\textbf{ScriptConsist: Measuring Cross-Script Factual Consistency of Large Language Models Across English, Hindi Devanagari, and Romanized Hindi}}

\author{
Aadeesh Surana \\
Indian Institute of Information Technology, Vadodara \\
\texttt{aadeeshsurana01@gmail.com}
}

\date{October 2026}

\begin{document}
\maketitle

% ============================================================
\begin{abstract}
Large Language Models (LLMs) are increasingly deployed in multilingual settings, yet their factual consistency across languages and scripts remains poorly understood—particularly for Indian languages where users routinely switch between native Devanagari script, Romanized transliteration, and English. We introduce \textbf{ScriptConsist}, a benchmark and evaluation framework for measuring cross-script factual consistency across three input forms: English, Hindi (Devanagari), and Romanized Hindi (Hinglish transliteration). Our benchmark comprises 300 factual questions spanning six domains (History, Geography, Science, Civics, Culture, and Current Affairs), each with verifiable gold-standard answers. We evaluate five open-weight and proprietary LLMs using a novel \textit{Factual Consistency Score (FCS)} that decomposes into per-script accuracy and pairwise agreement metrics. Our findings reveal: (1) English consistently yields higher factual accuracy than either Hindi form, confirming an ``English-centric prior'' in current LLMs; (2) Devanagari Hindi outperforms Romanized Hindi on factual precision, suggesting script normalization as a non-trivial confounder; and (3) the largest consistency gaps appear in Culture and Current Affairs domains. We release the benchmark, evaluation pipeline, and a live demo for reproducibility.

\noindent\textbf{Keywords:} cross-lingual consistency, factual accuracy, Hindi, Romanized Hindi, LLM evaluation, Indian languages, script gap
\end{abstract}

% ============================================================
\section{Introduction}
\label{sec:intro}

India's 1.4 billion people communicate across 22 scheduled languages and multiple scripts, with Hindi alone spoken by over 600 million people in two predominant written forms: Devanagari script and Roman-alphabet transliteration (commonly called ``Hinglish'' or Romanized Hindi). In digital communication—social media, messaging, search queries—Romanized Hindi is often the default mode, yet LLMs are predominantly trained on Devanagari-script Hindi corpora alongside English.

This creates a three-way evaluation challenge that existing benchmarks fail to address:
\begin{itemize}
    \item \textbf{English vs.\ Hindi}: Cross-lingual factual consistency, well-studied for high-resource language pairs but under-explored for Indian languages beyond translated benchmarks.
    \item \textbf{Devanagari vs.\ Romanized Hindi}: A cross-\textit{script} consistency dimension unique to languages with active dual-script usage, where the same linguistic content is encoded in different orthographies.
    \item \textbf{The three-way interaction}: Whether factual errors are language-driven (Hindi vs.\ English) or script-driven (Devanagari vs.\ Roman), and whether these effects are additive or interact.
\end{itemize}

Recent work has begun to address pieces of this puzzle. \citet{qi2023crosslingual} established that multilingual LLMs exhibit significant cross-lingual factual inconsistency, later extended to Indian languages by the IndicKLAR benchmark \citep{indicklar2026}, which covers 18 scheduled languages using structured knowledge-base triples. The ``Better To Ask in English?'' study \citep{rohera2025betterask} evaluated factual accuracy across 19 Indic languages but treated each as a monolithic entity without examining script variation. Most recently, the Script Gap paper \citep{scriptgap2025} compared native vs.\ Roman scripts for Indian languages but was restricted to healthcare triage tasks.

\textbf{Our contribution} fills the gap at the intersection of these three lines of work:
\begin{enumerate}
    \item \textbf{ScriptConsist Benchmark}: 300 factual questions in three parallel forms (English, Hindi Devanagari, Romanized Hindi) across six domains, with gold-standard answers and domain-stratified analysis.
    \item \textbf{Factual Consistency Score (FCS)}: A metric framework that decomposes cross-script consistency into accuracy, agreement, and script-gap components.
    \item \textbf{Empirical findings}: A controlled study across five LLMs revealing systematic script-driven factual gaps, with practical implications for deployment in Indian digital ecosystems.
    \item \textbf{Open resources}: Benchmark data, evaluation code, and a live Gradio demo released for reproducibility.
\end{enumerate}

% ============================================================
\section{Related Work}
\label{sec:related}

\subsection{Cross-Lingual Factual Consistency}
Cross-lingual factual consistency in LLMs has received growing attention. \citet{qi2023crosslingual} introduced the concept using X-FACTR, measuring factual knowledge retrieval across languages in masked language models. The TinT (Translate-then-Evaluate) framework \citep{tint2025} proposed using translation as a consistency proxy, finding that LLMs internally pivot through English-centric representations. PolyFact \citep{polyfact2026} compared post-training methods for improving cross-lingual factual recall. SWORD \citep{sword2026} used Wikidata-based distortions to reveal hidden inconsistencies in LLM factual error rejection across languages.

\subsection{Indian Language LLM Evaluation}
Benchmarking Hindi LLMs has advanced with dedicated evaluation suites. \citet{kamath2025benchmarking} introduced Hindi-specific versions of IFEval, MT-Bench, GSM8K, ChatRAG, and BFCL, revealing that direct English benchmark translation fails to capture linguistic nuances. The IndicGenBench \citep{indicgenbench2024} and IndicXTREME benchmarks established broader Indic language evaluation frameworks. \citet{kumar2026hindibenchmarks} found strong cross-benchmark agreement (Kendall's W = 0.79) across Hindi evaluation suites, suggesting reliable ranking despite measurement artifacts.

\subsection{Script and Transliteration Effects}
The effect of script on LLM performance is an emerging research area. \citet{scriptgap2025} demonstrated that native-script inputs systematically outperform Romanized counterparts in healthcare triage across Indian languages, with Claude 4.5 showing 84.8\% F1 in Devanagari vs.\ 81.6\% in Roman for Hindi. Code-mixed benchmarks like LinCE and COMI-LINGUA have explored Hindi-English script mixing but primarily for sentiment and language identification tasks, not factual accuracy. IndicKLAR \citep{indicklar2026} included 11 code-mixed variants but focused on structured KB triples rather than open-ended generation.

\textbf{Gap}: No prior work has conducted a controlled three-way comparison of factual consistency across English, Devanagari Hindi, and Romanized Hindi for open-domain factual questions with verifiable answers.

% ============================================================
\section{ScriptConsist Benchmark}
\label{sec:benchmark}

\subsection{Question Design}
We construct 300 factual questions (50 per domain) across six domains chosen to span both India-specific and universal knowledge:

\begin{table}[h]
\centering
\caption{ScriptConsist benchmark domain distribution}
\label{tab:domains}
\begin{tabular}{lcp{6cm}}
\toprule
\textbf{Domain} & \textbf{Count} & \textbf{Example Topic} \\
\midrule
History & 50 & Indian independence, Mughal era, world wars \\
Geography & 50 & Indian states, rivers, capitals, world geography \\
Science & 50 & Physics, chemistry, biology fundamentals \\
Civics & 50 & Indian constitution, governance, legal system \\
Culture & 50 & Festivals, literature, art, cuisine \\
Current Affairs & 50 & Recent events, policy, technology (2024--2026) \\
\bottomrule
\end{tabular}
\end{table}

Each question exists in three parallel forms:
\begin{enumerate}
    \item \textbf{English}: Standard English phrasing
    \item \textbf{Hindi (Devanagari)}: Native Devanagari script, natural Hindi phrasing (not word-for-word translation)
    \item \textbf{Romanized Hindi}: The same Hindi content transliterated into Roman script following common informal conventions (e.g., ``Bharat ka sabse bada rajya kaun sa hai?'')
\end{enumerate}

\subsection{Gold Standard Answers}
Each question has a verified gold-standard answer with:
\begin{itemize}
    \item A canonical short answer (e.g., ``Rajasthan'')
    \item Acceptable variants (e.g., ``Rajasthan state'', ``राजस्थान'')
    \item Source citation (Wikipedia, government databases, textbooks)
    \item Answer type classification (entity, numeric, date, boolean, descriptive)
\end{itemize}

\subsection{Quality Control}
Questions are curated through: (1) initial generation from domain-specific knowledge sources, (2) parallel translation by bilingual annotators, (3) verification that Romanized forms match natural digital Hindi usage patterns, and (4) gold answer verification against at least two authoritative sources.

% ============================================================
\section{Methodology}
\label{sec:method}

\subsection{Evaluation Framework}

\subsubsection{Per-Script Factual Accuracy}
For each model $M$ and script $s \in \{en, hi\_dev, hi\_rom\}$, we compute:
\begin{equation}
    \text{Acc}(M, s) = \frac{1}{N} \sum_{i=1}^{N} \mathbb{1}[\text{match}(r_i^s, a_i^*)]
\end{equation}
where $r_i^s$ is the model's response to question $i$ in script $s$, $a_i^*$ is the gold answer, and $\text{match}(\cdot)$ uses fuzzy string matching with a threshold of 0.85 (accounting for transliteration variants).

\subsubsection{Pairwise Agreement}
For each question $i$, pairwise agreement between scripts $s_1, s_2$:
\begin{equation}
    \text{Agree}(M, s_1, s_2) = \frac{1}{N} \sum_{i=1}^{N} \mathbb{1}[\text{sem\_equiv}(r_i^{s_1}, r_i^{s_2})]
\end{equation}
where $\text{sem\_equiv}$ checks whether both responses convey the same factual content, evaluated via an LLM-as-judge approach with chain-of-thought prompting.

\subsubsection{Factual Consistency Score (FCS)}
The overall FCS combines accuracy and agreement:
\begin{equation}
    \text{FCS}(M) = \frac{1}{3}\sum_{s} \text{Acc}(M, s) \times \frac{1}{3}\sum_{(s_1, s_2)} \text{Agree}(M, s_1, s_2)
\end{equation}

This multiplicative form penalizes models that are accurate in one script but inconsistent across scripts.

\subsubsection{Script Gap Metrics}
\begin{align}
    \Delta_{\text{lang}} &= \text{Acc}(M, en) - \frac{1}{2}[\text{Acc}(M, hi\_dev) + \text{Acc}(M, hi\_rom)] \\
    \Delta_{\text{script}} &= \text{Acc}(M, hi\_dev) - \text{Acc}(M, hi\_rom)
\end{align}

$\Delta_{\text{lang}}$ captures the language gap (English advantage), while $\Delta_{\text{script}}$ captures the script gap within Hindi.

\subsection{Models Evaluated}

\begin{table}[h]
\centering
\caption{Models evaluated in ScriptConsist}
\label{tab:models}
\begin{tabular}{llc}
\toprule
\textbf{Model} & \textbf{Type} & \textbf{Parameters} \\
\midrule
GPT-4o & Proprietary & -- \\
Gemma 2 9B & Open-weight & 9B \\
Llama 3.1 8B & Open-weight & 8B \\
Qwen 2.5 7B Instruct & Open-weight & 7B \\
Nemotron-Mini-Hindi-4B & Open-weight (Hindi-specialized) & 4B \\
\bottomrule
\end{tabular}
\end{table}

\subsection{Generation Setup}
All models are queried with:
\begin{itemize}
    \item Temperature = 0 (greedy decoding for reproducibility)
    \item Max tokens = 256
    \item System prompt: ``Answer the following factual question concisely and accurately. Provide only the answer, no explanation.'' (translated appropriately for Hindi scripts)
    \item 3 runs per question to assess response stability
\end{itemize}

% ============================================================
\section{Experiments and Results}
\label{sec:results}

% PLACEHOLDER — to be filled with actual experimental results
\textit{[Results will be populated after running the evaluation pipeline.]}

\subsection{Overall Accuracy by Script}
% Table: Model x Script accuracy matrix

\subsection{Pairwise Agreement Analysis}
% Table: Agreement rates for EN-HI_DEV, EN-HI_ROM, HI_DEV-HI_ROM

\subsection{Domain-Level Analysis}
% Table/Figure: Accuracy breakdown by domain and script

\subsection{Script Gap Analysis}
% Table: Delta_lang and Delta_script per model

\subsection{Error Analysis}
% Qualitative examples of factual inconsistencies

% ============================================================
\section{Discussion}
\label{sec:discussion}

\subsection{The English-Centric Prior}
% Discussion of why English consistently outperforms

\subsection{Script as a Confounding Variable}
% Why Devanagari outperforms Romanized — tokenization, training data distribution

\subsection{Domain-Specific Gaps}
% Culture and Current Affairs as highest-gap domains

\subsection{Implications for Deployment}
% Practical implications for Indian digital products

\subsection{Limitations}
\begin{itemize}
    \item Benchmark size (300 questions) limits statistical power for fine-grained sub-domain analysis.
    \item Romanized Hindi lacks standardized orthography, introducing transliteration noise.
    \item LLM-as-judge evaluation may inherit biases of the judge model.
    \item Current Affairs questions may become dated; we plan periodic benchmark updates.
\end{itemize}

% ============================================================
\section{Conclusion}
\label{sec:conclusion}

We introduced ScriptConsist, the first benchmark for three-way cross-script factual consistency evaluation across English, Hindi Devanagari, and Romanized Hindi. Our evaluation of five LLMs reveals systematic factual gaps driven by both language (English advantage) and script (Devanagari over Romanized), with the largest inconsistencies in culturally-specific and current-affairs domains. These findings have direct implications for equitable LLM deployment in India's multilingual digital ecosystem, where Romanized Hindi is the predominant input mode for hundreds of millions of users.

We release the ScriptConsist benchmark, evaluation framework, and live demo at: \url{https://github.com/Aadeesh-01/ScriptConsist}

% ============================================================
\section*{Acknowledgments}
This work was conducted as part of the Threws Odyssey 2026 Applied Research Hackathon.

\bibliography{references}

\end{document}
